% ==============================================================================
% MÓDULO CAPÍTULO II: CARACTERÍSTICAS SOCIOCULTURALES, LINGÜÍSTICAS Y COSMOVISIÓN
% Archivo: cap_2/2.3_sociocultural_quechua.tex
% ==============================================================================
\subsection{Características socioculturales, lingüísticas y cosmovisión andina}
\label{sec:2.3_social}

El entorno comunitario de Belén se distingue por una singular riqueza identitaria y una densa trama de sociabilidad popular. Históricamente, el barrio de Belén ha sido cuna de ilustres artesanos populares de Huamanga, reconocidos internacionalmente por el arte del retablo ayacuchano, la imaginería religiosa en yeso y el tallado tradicional en piedra de Huamanga. En este espacio geográfico coexisten familias tradicionales residentes desde hace generaciones con oleadas de población migrante provenientes de las zonas rurales y provincias del centro-sur ayacuchano (Víctor Fajardo, Cangallo y Vilcas Huamán), muchas de ellas asentadas a raíz de la violencia sociopolítica de las décadas pasadas.

En el plano lingüístico, el \textbf{72\% de la población usuaria es bilingüe activa en quechua chanka y castellano}. En las generaciones de adultos mayores y en mujeres migrantes de sectores periféricos, el quechua constituye la lengua materna predominante y el vehículo exclusivo para verbalizar el sufrimiento corpóreo y anímico. Para la cosmovisión andina huamanguina, los padecimientos no son meras disfunciones biológicas aisladas, sino alteraciones del equilibrio integral del ser humano con su entorno, expresadas en nociones complejas como el dolor visceral (\textit{nanay}), el estado de enfermedad integral (\textit{onqoy}), la tristeza y aflicción acumulada (\textit{llakikuy}), el enfriamiento del cuerpo (\textit{chiriyay}) o los desórdenes espirituales derivados del espanto o susto (\textit{mancharisqa}).

Las pautas de sociabilidad comunitaria exigen de manera indeclinable la observancia de códigos éticos y de cortesía tradicionales basados en el respeto mutuo (\textit{respetanakuy}), el saludo formal y afectuoso (\textit{napaykuy}) y la acogida hospitalaria (\textit{allin chaskiy}) \cite{aguirre2020}. Cuando el personal de salud —frecuentemente hispanohablante monolingüe o condicionado por la prisa protocolar— prescinde del saludo, evita el contacto visual, interrumpe el relato del usuario o emplea un léxico excesivamente tecnicista, se produce un violento quiebre comunicativo. El usuario andino interpreta esta conducta como desdén, frialdad burocrática, soberbia institucional o discriminación clasista y étnica, replegándose en el silencio (\textit{upallay}) y perdiendo la confianza en el sistema asistencial.

En cuanto al nivel de instrucción, el 48\% de los jefes de familia cuenta con educación secundaria completa, un 26\% posee secundaria incompleta o primaria, y un 8\% (concentrado en adultos mayores) registra analfabetismo formal o funcional. Esta realidad representa un desafío crítico para la alfabetización en salud, pues muchos usuarios no logran descifrar las recetas médicas escritas, los calendarios de citas o la señalética institucional de los consultorios, dependiendo enteramente de la orientación verbal empática que el profesional de enfermería pueda suministrarle.
